\chapter{Nota final: lo real y lo inventado}
\markboth{Lo real y lo inventado}{Lo real y lo inventado}

\inicio{T}{odo} lo que ocurre en Ceres es inventado: el Archivo, la Bóveda, la \emph{Perséfone}, Inés Calderón, Siete y el resto de los personajes. Las Tres Reglas del Archivo son un homenaje a la manera en que los viejos maestros de la ciencia ficción convertían un problema técnico en un dilema entre leyes; no aparecen en ningún paper.

Lo que no es inventado:

\begin{itemize}
\item \textbf{El problema.} La carrera \emph{read-reclaim} existe en el conteo de referencias concurrente, y el \emph{use-after-free} que provoca es una causa habitual de vulnerabilidades.
\item \textbf{La solución.} El talonario es el \emph{peso} del paper de de Wolff, Anderson, Keller y Seletskiy (Euro-Par 2024). Los procedimientos de los capítulos IV a VI son, paso a paso, sus algoritmos 1 y 2. Los requisitos R-1 y R-2, la ofensa, el multiconjunto $O$ y las dos afirmaciones del capítulo VII son su demostración (secciones 4 y 4.1).
\item \textbf{Las extensiones y los límites} del capítulo VIII están en la sección 5 del paper (punteros nulos, débiles, \emph{compare-and-swap} y punteros \emph{aliased}) o se derivan de él. Las observaciones sobre ARM, LA57, el orden de memoria y la ausencia de \emph{snapshot pointers} proceden en parte de la explicación detallada en la que se basa este libro.
\item \textbf{Los ensayos del capítulo X} reproducen los resultados de la sección 6 del paper: pilas concurrentes en una máquina de 144 hilos y latencias en una de 56. Los valores de la tabla de pilas son lecturas aproximadas de su Figura 5; los de latencia, las medianas de su Figura 6.
\item \textbf{El teatro y el ensayo de estrés del capítulo IX} no están en el paper: corresponden al simulador y al boceto en C++ de la explicación detallada, con sus resultados.
\item \textbf{El trabajo previo} que recita el profesor Loew en el capítulo III es real: Valois; Michael y Scott; Detlefs, Martin, Moir y Steele; Herlihy, Luchangco, Martin y Moir; Gidenstam et al.; Sundell; Anderson, Blelloch y Wei; el \emph{split reference count} de Lee, Williams y Folly; Bevan y Watson y Watson para el conteo ponderado.
\item \textbf{Y sí:} Daniel Anderson es a la vez coautor del paper y autor de la implementación rival que lleva su nombre en los experimentos.
\end{itemize}

\chapter{Apéndice A · Diccionario de Ceres}
\markboth{Diccionario de Ceres}{Diccionario de Ceres}

\noindent Cada objeto del Archivo, con su nombre en el paper y lo que es de verdad.

\medskip
{\small
\begin{longtable}{@{}>{\itshape\raggedright}p{2.5cm} >{\raggedright}p{2.7cm} >{\raggedright\arraybackslash}p{5.4cm}@{}}
\toprule
\upshape\textbf{En Ceres} & \textbf{En el paper} & \textbf{Qué es} \\
\midrule\endhead
Bóveda & memoria (\emph{heap}) & donde viven los objetos \\
nicho & dirección & número de una posición de memoria \\
cristal & objeto & datos con contador de referencias \\
hoja de usuarios & \emph{reference count} (rc) & entero junto al objeto \\
fundidor & liberar (\emph{reclaim}) & devolver la memoria; puede reutilizarse \\
cartel & \emph{atomic shared pointer} & celda mutable compartida \\
ticket, pase & \emph{shared pointer} (peso 1) & referencia propia e inmutable \\
talonario & peso (\emph{weight}) & bits altos del puntero empaquetado \\
lector & hilo que hace \texttt{load} & hasta $P$ hilos \\
archivero & hilo que hace \texttt{store} & escritor \\
censo, $P$ & $P$ & número máximo de hilos \\
talonario lleno, $C$ & $C$ & peso con que se carga; $C\geq 2P$ \\
arrancar & \texttt{fetch\_and\_add} & suma y devuelve el valor anterior \\
trocar & \texttt{exchange} / \emph{swap} & escribe y devuelve el anterior \\
cotejar & \texttt{compare\_and\_swap} & cambia solo si coincide \\
mano de dos brazos & DCAS & CAS sobre dos palabras separadas; no existe \\
placa doble & CAS de 128 bits & sí existe, pero sin \texttt{fetch\_and\_add} \\
zona roja & peso $\leq P$ & hay que recargar \\
recargar & \texttt{increase\_weight} & algoritmo 2 \\
ofensa & \emph{offence} & decrementar con peso $\leq P$ \\
salida temprana & \emph{break} si el peso subió & lo que lo hace \emph{wait-free} \\
llave & candado (\emph{lock}) & la versión de libstdc++ \\
basura diferida & SMR, \emph{hazard pointers}, épocas & reclamación diferida \\
vistazo & \emph{snapshot pointer} & lectura sin tocar el contador \\
contador partido & \emph{split reference count} & peso $= C -$ contador local \\
cita & puntero \emph{aliased} & comparte contador, apunta a otra cosa \\
copista en el sitio & actualización destructiva & reutilizar si rc $=1$ \\
cartel vacío & \texttt{null} & puntero nulo \\
ficha débil & \emph{weak pointer} & referencia que no mantiene vivo \\
Cerebro & la máquina de pruebas & 144 hilos (rendimiento), 56 (latencia) \\
\bottomrule
\end{longtable}
}

\chapter{Apéndice B · Los algoritmos, en el idioma de Inés}
\markboth{Los algoritmos}{Los algoritmos}

\noindent Versión en español de los algoritmos 1 y 2 del paper. \textit{empaquetar}$(p,w)$ junta nicho y talonario en una placa; \textit{desempaquetar} los separa. Un ticket vale $2^{48}$ porque el talonario vive en las casillas altas.

\begin{pizarra}
\ttfamily\footnotesize
\textbf{cambiar}(cartel, nuevo):\\
\hspace*{1.2em}hoja(nuevo) += C\\
\hspace*{1.2em}v ← trocar(cartel, empaquetar(nuevo, C))\\
\hspace*{1.2em}(viejo, w) ← desempaquetar(v)\\
\hspace*{1.2em}hoja(viejo) -= w \quad\textrm{\itshape // si llega a 0, al fundidor}\\[2mm]
\textbf{leer}(cartel):\\
\hspace*{1.2em}v ← arrancar(cartel, -empaquetar(nulo, 1))\\
\hspace*{1.2em}(p, w) ← desempaquetar(v) \quad\textrm{\itshape // v = valor anterior}\\
\hspace*{1.2em}si w - 1 < P: recargar(cartel, p, w - 1)\\
\hspace*{1.2em}devolver ticket(p) \quad\textrm{\itshape // ya apuntado en la hoja}
\end{pizarra}

\begin{pizarra}
\ttfamily\footnotesize
\textbf{recargar}(cartel, p, w):\\
\hspace*{1.2em}hoja(p) += C - w \quad\textrm{\itshape // reservar ANTES de cotejar}\\
\hspace*{1.2em}esperado ← w\\
\hspace*{1.2em}repetir:\\
\hspace*{2.4em}r ← cotejar(cartel, empaquetar(p, esperado),\\
\hspace*{8.4em}empaquetar(p, esperado + C - w))\\
\hspace*{2.4em}si r.éxito: terminar\\
\hspace*{2.4em}(p2, w2) ← desempaquetar(r.visto)\\
\hspace*{2.4em}si p2 ≠ p o w2 > esperado: \textrm{\itshape // salida temprana}\\
\hspace*{3.5999999999999996em}salir del bucle\\
\hspace*{2.4em}esperado ← w2\\
\hspace*{1.2em}hoja(p) -= C - w \quad\textrm{\itshape // devolver la reserva}
\end{pizarra}

\noindent\textbf{Cambiar con condiciones} (\emph{compare-and-set}, algoritmo 3, solo \emph{lock-free}): apuntar $C$ en el nuevo; leer el cartel; mientras el cristal visto sea el esperado, cotejar «(visto, tickets vistos) $\to$ (nuevo, $C$)»; si acepta, restar en el viejo los tickets vistos y responder sí; si el cristal visto es otro, devolver los $C$ del nuevo y responder no; si solo cambiaron los tickets, reintentar.

\chapter{Apéndice C · Las cuentas de la pizarra}
\markboth{Las cuentas}{Las cuentas}

\begin{description}[leftmargin=0pt,itemsep=2mm,font=\normalfont\scshape\color{oxido}]
\item[La regla de la hoja] hoja $\geq$ talonario $+$ tickets fuera; igualdad cuando nadie está a mitad de un procedimiento.
\item[R-1] $C \geq 2P$.
\item[R-2] el talonario baja de uno en uno o salta por encima de $P$.
\item[Por qué la recarga sale de la zona roja] $\text{actual} + C - w \;\geq\; 0 + 2P - (P-1) \;=\; P+1 \;>\; P$.
\item[Afirmación 1] cada lector está en $O$ como mucho una vez $\Rightarrow |O| \leq P$.
\item[Afirmación 2] talonario $\geq P - |O|$.
\item[Conclusión] talonario $\geq P - P = 0$: nunca negativo.
\item[Intentos de recarga] los valores esperados forman una cadena estrictamente decreciente, por debajo de $P$ y $\geq 0$: como mucho $P$ intentos.
\item[Casillas] 16 libres $\Rightarrow C \leq 2^{16}-1 = 65\,535$; con 3 de alineación, 19 $\Rightarrow C \leq 2^{19}-1$. En los experimentos: $C = 2^{19}$, $P = 2^{16}$.
\item[Lecturas entre recargas] $C - P = 524\,288 - 65\,536 = 458\,752$: unas dos recargas por millón de lecturas.
\item[Hojas de 64 cifras] con hojas de 32 bits y $C=2^{19}$, bastan $2^{32}/2^{19} = 8\,192$ carteles para desbordar; con 64, harían falta $2^{45}$.
\item[Contador partido] peso $= C -$ contador local. Efecto neto en la hoja tras un cambio: $+C - (C - \text{contador}) = +\text{contador}$.
\item[Costes] leer: casi siempre 1 gesto, como mucho $O(P)$; cambiar y trocar: $O(1)$; cambiar con condiciones y citas: \emph{lock-free}.
\end{description}

\chapter{Apéndice D · Del paper a los capítulos}
\markboth{Del paper a los capítulos}{Del paper a los capítulos}

{\small
\begin{tabularx}{\textwidth}{@{}l X l@{}}
\toprule
\textbf{Sección del paper} & \textbf{Idea} & \textbf{Capítulo} \\
\midrule
Resumen y §1 & el problema y por qué importa & Prólogo, I, II \\
§2 Preliminares & interfaz; carreras \emph{read-reclaim}; SMR & II, III \\
§3 Conteo ponderado & pesos; algoritmo 1; empaquetado & IV \\
§4 Aumentar el peso & recarga; R-1, R-2; algoritmo 2 & V, VI \\
§4.1--4.2 & demostración y coste & VII \\
§5 Extensiones & nulos, débiles, CAS, \emph{aliasing} & VIII \\
§6 Experimentos & pilas y latencias & X \\
§7 Trabajo relacionado & historia; \emph{split reference count} & III, VII, Epílogo \\
§8 Conclusión & wait-free sin reclamación diferida & X, Epílogo \\
\bottomrule
\end{tabularx}
}

\medskip
\noindent\textbf{Referencia.} I.~G.~de Wolff, D.~Anderson, G.~K.~Keller, A.~Seletskiy. \emph{A Fast Wait-Free Solution to Read-Reclaim Races in Reference Counting}. Euro-Par 2024, LNCS 14803, pp.~103--118. Código y experimentos: Zenodo y GitHub (\texttt{ivogabe/wait-free-atomic-shared-pointer}).

\vfill
\begin{center}
{\color{oxido}$\ast$}\\[2mm]
{\small\itshape Fin}
\end{center}
